% ---------------------------------------------------------------------
%  3. Marco teórico — D (definiciones), A (Huffman), B (LZW),
%     C (comparación teórica e hipótesis)
% ---------------------------------------------------------------------
\section{Marco teórico}
\label{sec:marco-teorico}

\subsection{Definiciones preliminares}
\label{sec:definiciones}
\begin{aporte}{D}
Sea $x = x_1x_2\cdots x_n$ una secuencia de $n$ bytes sobre el alfabeto
$\Sigma = \{0,\dots,255\}$, de la cual aparecen $\sigma \le 256$ símbolos distintos. La
frecuencia del símbolo $s$ es $f_s$ y su probabilidad empírica es $p_s = f_s/n$.

\paragraph{Fuente y entropía de orden cero.} Un modelo de orden cero trata los símbolos como
independientes e idénticamente distribuidos (i.i.d.). Su entropía empírica
\begin{equation}
  H_0(x) = -\sum_{s\,:\,f_s>0} p_s \log_2 p_s \qquad \text{[bits por símbolo]}
  \label{eq:entropia}
\end{equation}
es el número mínimo promedio de bits por símbolo de cualquier código que asigne a cada
símbolo una palabra fija~\cite{shannon1948}. Cumple $0 \le H_0 \le \log_2\sigma \le 8$, con
igualdad superior para la distribución uniforme.

\paragraph{Códigos prefijo y desigualdad de Kraft--McMillan.} Un código asigna a cada símbolo
$s$ una palabra binaria de longitud $\ell_s$; es \emph{prefijo} si ninguna palabra es prefijo
de otra, lo que permite decodificar sin separadores. Todo código unívocamente decodificable
satisface
\begin{equation}
  \sum_{s} 2^{-\ell_s} \le 1,
  \label{eq:kraft}
\end{equation}
y, recíprocamente, para cualquier conjunto de longitudes que cumpla \eqref{eq:kraft} existe un
código prefijo con esas longitudes~\cite{sayood2017,salomon2010}. La longitud media del código
es $\Lbar = \sum_s p_s \ell_s$.

\paragraph{Métricas de compresión.} Si $|C(x)|$ es el tamaño en bytes de la salida comprimida,
incluida la cabecera, se definen
\begin{equation}
  \bpb = \frac{8\,|C(x)|}{n}, \qquad
  \text{razón} = \frac{n}{|C(x)|}, \qquad
  \text{ahorro} = 1 - \frac{|C(x)|}{n}.
  \label{eq:metricas}
\end{equation}
Los bits por byte (bpb) permiten comparar directamente con $H_0$: un valor $\bpb > 8$
indica expansión.
\end{aporte}

% =====================================================================
\subsection{Codificación de Huffman}
\label{sec:huffman}
\begin{aporte}{A}
\paragraph{Idea.} Huffman~\cite{huffman1952} construye un árbol binario de abajo hacia arriba:
parte de una hoja por símbolo, ponderada por su frecuencia, y fusiona repetidamente los dos
árboles de menor peso. La profundidad de cada hoja es la longitud de su código. Se trata de un
algoritmo voraz cuya elección local ---fusionar los dos símbolos menos frecuentes--- es
óptima~\cite{cormen2022}. El \cref{alg:huffman} muestra el cálculo de longitudes tal como lo
realiza \texttt{HuffmanCodec.codeLengths}, con un montículo binario de mínimos y desempate
por orden de creación del nodo para que el resultado sea determinista.

\begin{algorithm}[H]
  \caption{Longitudes de código de Huffman (\texttt{codeLengths}).}
  \label{alg:huffman}
  \small
  \begin{algorithmic}[1]
    \Require frecuencias $f_0,\dots,f_{255}$
    \Ensure longitudes $\ell_0,\dots,\ell_{255}$ ($\ell_s = 0$ si $f_s = 0$)
    \State $Q \gets$ montículo de mínimos con una hoja por cada $s$ con $f_s > 0$, clave $w(s)=f_s$
    \If{$|Q| = 1$} \State \Return $\ell_s = 1$ para el único símbolo presente \EndIf
    \While{$|Q| > 1$}
      \State $a \gets \textsc{ExtraerMín}(Q)$;\; $b \gets \textsc{ExtraerMín}(Q)$
      \State crear nodo $z$ con $w(z) = w(a) + w(b)$;\; $\mathit{padre}(a) \gets z$;\; $\mathit{padre}(b) \gets z$
      \State $\textsc{Insertar}(Q, z)$
    \EndWhile
    \ForAll{nodo $v$ en orden decreciente de creación}
      \State $\mathit{prof}(v) \gets \mathit{prof}(\mathit{padre}(v)) + 1$ \Comment{la raíz tiene profundidad 0}
    \EndFor
    \State \Return $\ell_s = \mathit{prof}(s)$ para cada hoja $s$
  \end{algorithmic}
\end{algorithm}

\paragraph{Optimalidad y cota de entropía.} Entre todos los códigos prefijo para una
distribución dada, el de Huffman minimiza $\Lbar$~\cite{huffman1952,cormen2022}. Como su árbol
es completo (todo nodo interno tiene dos hijos), para $\sigma \ge 2$ la desigualdad
\eqref{eq:kraft} se cumple con igualdad. Además,
\begin{equation}
  H_0 \;\le\; \Lbar \;<\; H_0 + 1,
  \label{eq:cota-huffman}
\end{equation}
la cota inferior por el teorema de Shannon~\cite{shannon1948} y la superior porque el código
con longitudes $\lceil -\log_2 p_s\rceil$ satisface \eqref{eq:kraft} y Huffman no puede ser
peor~\cite{sayood2017}. Gallager~\cite{gallager1978} refinó la cota superior en función de la
probabilidad del símbolo más frecuente.

\paragraph{Códigos canónicos.} Las longitudes bastan para reconstruir el código si las palabras
se asignan de forma \emph{canónica}: en orden creciente de longitud y, dentro de cada longitud,
en orden creciente de símbolo, como especifica DEFLATE~\cite{rfc1951}. Con $c_\ell$ el número
de códigos de longitud $\ell$, el primer código de longitud $\ell$ es
\begin{equation}
  \mathit{first}_1 = 0, \qquad \mathit{first}_\ell = \left(\mathit{first}_{\ell-1} + c_{\ell-1}\right) \cdot 2,
  \label{eq:canonico}
\end{equation}
y los símbolos de esa longitud reciben códigos consecutivos a partir de él. Por eso la
implementación guarda en la cabecera solo 256 longitudes de un byte, en lugar del árbol.

\paragraph{Ejemplo: \texttt{abracadabra}.} La cadena tiene $n = 11$ y frecuencias
$a{:}5$, $b{:}2$, $r{:}2$, $c{:}1$, $d{:}1$. El algoritmo fusiona primero $c$ y $d$ (peso 2),
luego $b$ y $r$ (peso 4), después ambos subárboles (peso 6) y finalmente el resultado con $a$
(peso 11). La \cref{fig:arbol-huffman} muestra el árbol y la \cref{tab:abracadabra} las
longitudes y los códigos. El código canónico depende solo de las longitudes, por lo que puede
diferir del que se lee en el árbol sin cambiar $\Lbar$.

\begin{figure}[H]
  \centering
  \begin{tikzpicture}[
      level distance=1.25cm,
      level 1/.style={sibling distance=4.6cm},
      level 2/.style={sibling distance=2.8cm},
      level 3/.style={sibling distance=1.4cm},
      interno/.style={circle, draw=black!70, fill=gray!10, minimum size=8mm, inner sep=0pt, font=\small},
      hoja/.style={rectangle, rounded corners=2pt, draw=black!70, fill=white, minimum width=11mm,
                   minimum height=7mm, font=\small},
      edge from parent/.style={draw=black!55, -{Stealth[length=2mm]}},
      etiqueta/.style={font=\scriptsize, fill=white, inner sep=1pt}]
    \node[interno] {11}
      child { node[hoja] {\texttt{a}\,:\,5} edge from parent node[etiqueta, above left] {0} }
      child { node[interno] {6}
        child { node[interno] {2}
          child { node[hoja] {\texttt{c}\,:\,1} edge from parent node[etiqueta, above left] {0} }
          child { node[hoja] {\texttt{d}\,:\,1} edge from parent node[etiqueta, above right] {1} }
          edge from parent node[etiqueta, above left] {0} }
        child { node[interno] {4}
          child { node[hoja] {\texttt{b}\,:\,2} edge from parent node[etiqueta, above left] {0} }
          child { node[hoja] {\texttt{r}\,:\,2} edge from parent node[etiqueta, above right] {1} }
          edge from parent node[etiqueta, above right] {1} }
        edge from parent node[etiqueta, above right] {1} };
  \end{tikzpicture}
  \caption{Árbol de Huffman para \texttt{abracadabra}; los nodos internos muestran el peso acumulado.}
  \label{fig:arbol-huffman}
\end{figure}

\begin{table}[H]
  \centering
  \caption{Longitudes y códigos para \texttt{abracadabra}.}
  \label{tab:abracadabra}
  \small
  \begin{tabular}{@{}c c c c c@{}}
    \toprule
    Símbolo & Frecuencia & Longitud $\ell_s$ & Código del árbol & Código canónico \\
    \midrule
    \texttt{a} & 5 & 1 & \texttt{0}   & \texttt{0}   \\
    \texttt{b} & 2 & 3 & \texttt{110} & \texttt{100} \\
    \texttt{c} & 1 & 3 & \texttt{100} & \texttt{101} \\
    \texttt{d} & 1 & 3 & \texttt{101} & \texttt{110} \\
    \texttt{r} & 2 & 3 & \texttt{111} & \texttt{111} \\
    \bottomrule
  \end{tabular}
\end{table}

Con $c_1 = 1$, $c_2 = 0$ y $c_3 = 4$, la \cref{eq:canonico} da $\mathit{first}_1 = 0$,
$\mathit{first}_2 = (0 + 1)\cdot 2 = 2$ y $\mathit{first}_3 = (2 + 0)\cdot 2 = 4 = \texttt{100}_2$, de modo que $b, c, d, r$
reciben \texttt{100}, \texttt{101}, \texttt{110} y \texttt{111}. La cadena se codifica como
\texttt{0\,100\,111\,0\,101\,0\,110\,0\,100\,111\,0}, es decir, 23 bits frente a 88 bits sin
comprimir. Así, $\Lbar = 23/11 \approx 2.091$ y $H_0 \approx 2.040$, lo que verifica
\eqref{eq:cota-huffman}; además $2^{-1} + 4\cdot 2^{-3} = 1$, la igualdad de Kraft.

\paragraph{Decodificación canónica.} Siguiendo a Moffat y Turpin~\cite{moffat1997}, el
decodificador no reconstruye el árbol: lee un bit por iteración y, para cada longitud $\ell$,
comprueba si el valor acumulado cae en el rango de los $c_\ell$ códigos de esa longitud
(\cref{alg:huffman-dec}). Los símbolos se almacenan ordenados por código canónico, de modo que
el índice del símbolo se obtiene por simple resta.

\begin{algorithm}[H]
  \caption{Decodificación canónica de un símbolo (\texttt{decodeSymbol}).}
  \label{alg:huffman-dec}
  \small
  \begin{algorithmic}[1]
    \Require lector de bits; $c_\ell$ por longitud; símbolos ordenados por código canónico
    \State $\mathit{code} \gets 0$;\; $\mathit{first} \gets 0$;\; $\mathit{index} \gets 0$
    \For{$\ell = 1$ \textbf{hasta} $L_{\max}$}
      \State $\mathit{code} \gets \mathit{code} \mathbin{|} \textsc{LeerBit}()$
      \If{$\mathit{code} - \mathit{first} < c_\ell$}
        \State \Return $\textit{símbolos}[\mathit{index} + \mathit{code} - \mathit{first}]$
      \EndIf
      \State $\mathit{index} \gets \mathit{index} + c_\ell$;\;
             $\mathit{first} \gets (\mathit{first} + c_\ell)\cdot 2$;\;
             $\mathit{code} \gets \mathit{code}\cdot 2$
    \EndFor
    \State \textbf{error}: código inválido
  \end{algorithmic}
\end{algorithm}

\paragraph{Análisis de complejidad.} Sea $n$ el tamaño de la entrada y $\sigma \le 256$ el
número de símbolos presentes.
\begin{itemize}
  \item \emph{Conteo de frecuencias:} una pasada con un incremento por byte, $\Theta(n)$.
  \item \emph{Construcción del código:} $\sigma$ inserciones iniciales y $\sigma - 1$
        iteraciones con dos extracciones y una inserción en el montículo, cada una
        $O(\log\sigma)$; en total $O(\sigma\log\sigma)$. La asignación de profundidades y de
        códigos canónicos cuesta $O(\sigma + L_{\max})$.
  \item \emph{Codificación:} segunda pasada sobre la entrada; cada símbolo escribe $\ell_s$
        bits con \texttt{BitWriter} en tiempo $O(1 + \ell_s/8)$, de modo que el total es
        $\Theta(n)$. El costo global de codificar es $\Theta(n + \sigma\log\sigma)$, es decir,
        $\Theta(n)$ para un alfabeto fijo.
  \item \emph{Decodificación:} el \cref{alg:huffman-dec} realiza $\ell_s$ iteraciones por
        símbolo, por lo que el total es $\sum_i \ell_{x_i} = n\Lbar$ iteraciones, es decir,
        $\Theta(n\Lbar)$. Por \eqref{eq:cota-huffman}, $\Lbar < H_0 + 1 \le \log_2 256 + 1 = 9$,
        así que el costo es lineal en $n$, pero con una constante proporcional a la longitud
        media del código y, por tanto, a la entropía de la fuente.
  \item \emph{Espacio adicional:} $O(\sigma)$ para el montículo, los pesos, los padres y las
        tablas canónicas, sin contar los arreglos de entrada y salida. La cabecera ocupa
        $2 + 4 + 256 = 262$ bytes, un sobrecosto de $8\cdot 262/n$ bpb que solo es apreciable
        en entradas pequeñas.
\end{itemize}
\end{aporte}

% =====================================================================
\subsection{Algoritmo Lempel--Ziv--Welch (LZW)}
\label{sec:lzw}
\begin{aporte}{B}
\paragraph{Idea.} LZW~\cite{welch1984} es una variante de LZ78~\cite{ziv1978} que inicializa el
diccionario con los 256 bytes posibles. El codificador extiende la frase actual $w$ mientras
$w\,c$ siga en el diccionario; cuando no lo está, emite el código de $w$, agrega $w\,c$ como
nueva entrada y reinicia la frase con $c$. El decodificador reconstruye el mismo diccionario
a partir de los códigos, sin que este se transmita. En esta implementación el código 256 se
reserva para \texttt{CLEAR}, de modo que las frases nuevas comienzan en 257.

\paragraph{Códigos de ancho variable y reinicio.} Los códigos se escriben con el menor ancho
capaz de representar todos los códigos posibles en ese momento,
$\mathit{ancho}(s) = \max\left(9, \lceil\log_2 s\rceil\right)$, donde $s$ es el siguiente
código libre; el ancho crece de 9 hasta $b$ bits, con $b \in [9, 16]$ configurable
(\texttt{maxBits}, 16 por defecto). Cuando el diccionario alcanza $2^b$ entradas, el
codificador emite \texttt{CLEAR} y ambos lados vuelven al diccionario inicial. Esta política
de reinicio mantiene el modelo adaptativo ante cambios en los datos y acota la memoria, a costa
de olvidar frases útiles~\cite{cheng2021}. Los \cref{alg:lzw-enc,alg:lzw-dec} resumen
\texttt{LzwCodec.compress} y \texttt{LzwCodec.decode}.

\begin{algorithm}[H]
  \caption{Codificación LZW con reinicio (\texttt{compress}).}
  \label{alg:lzw-enc}
  \small
  \begin{algorithmic}[1]
    \Require bytes $x_1\cdots x_n$ ($n \ge 1$), ancho máximo $b$
    \State $D \gets$ diccionario con los 256 bytes;\; $s \gets 257$;\; $w \gets x_1$
    \For{$i = 2$ \textbf{hasta} $n$}
      \State $c \gets x_i$
      \If{$(w, c) \in D$}
        \State $w \gets D[w, c]$ \Comment{extender la frase}
      \Else
        \State emitir $w$ con $\mathit{ancho}(s)$ bits
        \If{$s < 2^b$}
          \State $D[w, c] \gets s$;\; $s \gets s + 1$
        \Else
          \State emitir \texttt{CLEAR} con $\mathit{ancho}(s)$ bits;\; vaciar $D$;\; $s \gets 257$
        \EndIf
        \State $w \gets c$
      \EndIf
    \EndFor
    \State emitir $w$ con $\mathit{ancho}(s)$ bits
  \end{algorithmic}
\end{algorithm}

\begin{algorithm}[H]
  \caption{Decodificación LZW (\texttt{decode}).}
  \label{alg:lzw-dec}
  \small
  \begin{algorithmic}[1]
    \Require flujo de códigos, longitud original $n$, ancho máximo $b$
    \State $s \gets 257$;\; $\mathit{prev} \gets \bot$
    \While{se hayan producido menos de $n$ bytes}
      \State leer $k$ con $\mathit{ancho}(s)$ bits si $\mathit{prev} = \bot$, o con
             $\mathit{ancho}(\min(s + 1, 2^b))$ en otro caso
      \If{$k = \texttt{CLEAR}$} \State $s \gets 257$;\; $\mathit{prev} \gets \bot$;\; \textbf{continuar} \EndIf
      \If{$\mathit{prev} = \bot$} \State producir el byte $k$;\; $\mathit{prev} \gets k$;\; \textbf{continuar} \EndIf
      \If{$s < 2^b$}
        \State $h \gets \mathit{primero}(k)$ si $k < s$; si no, $h \gets \mathit{primero}(\mathit{prev})$ \Comment{caso $KwKwK$}
        \State $D[s] \gets (\mathit{prev}, h)$;\; $s \gets s + 1$
      \EndIf
      \State producir la frase $D[k]$ recorriendo prefijos de atrás hacia adelante;\; $\mathit{prev} \gets k$
    \EndWhile
  \end{algorithmic}
\end{algorithm}

El decodificador agrega cada entrada un código más tarde que el codificador, porque necesita
el primer byte de la frase siguiente; por eso lee con el ancho correspondiente a $s + 1$.

\paragraph{Ejemplo: \texttt{TOBEORNOTTOBEORTOBEORNOT}.} La \cref{tab:lzw-welch} reproduce el
ejemplo clásico de Welch con la numeración de esta implementación. Las 24 letras se codifican
con 16 códigos, la secuencia
\texttt{T O B E O R N O T 257 259 261 266 260 262 264}, que es exactamente la que verifica la
prueba \texttt{emitsTextbookCodeSequence}. Como el diccionario nunca supera 512 entradas, todos
los códigos usan 9 bits: $16 \cdot 9 = 144$ bits frente a $24 \cdot 8 = 192$ bits de entrada.

\begin{table}[H]
  \centering
  \caption{Codificación LZW de \texttt{TOBEORNOTTOBEORTOBEORNOT} (solo los pasos que emiten).}
  \label{tab:lzw-welch}
  \small
  \begin{tabular}{@{}r l c r l@{}}
    \toprule
    Paso & Frase $w$ & Siguiente $c$ & Código emitido & Nueva entrada \\
    \midrule
     1 & \texttt{T}   & \texttt{O} & 84 (\texttt{T}) & 257 = \texttt{TO} \\
     2 & \texttt{O}   & \texttt{B} & 79 (\texttt{O}) & 258 = \texttt{OB} \\
     3 & \texttt{B}   & \texttt{E} & 66 (\texttt{B}) & 259 = \texttt{BE} \\
     4 & \texttt{E}   & \texttt{O} & 69 (\texttt{E}) & 260 = \texttt{EO} \\
     5 & \texttt{O}   & \texttt{R} & 79 (\texttt{O}) & 261 = \texttt{OR} \\
     6 & \texttt{R}   & \texttt{N} & 82 (\texttt{R}) & 262 = \texttt{RN} \\
     7 & \texttt{N}   & \texttt{O} & 78 (\texttt{N}) & 263 = \texttt{NO} \\
     8 & \texttt{O}   & \texttt{T} & 79 (\texttt{O}) & 264 = \texttt{OT} \\
     9 & \texttt{T}   & \texttt{T} & 84 (\texttt{T}) & 265 = \texttt{TT} \\
    10 & \texttt{TO}  & \texttt{B} & 257 & 266 = \texttt{TOB} \\
    11 & \texttt{BE}  & \texttt{O} & 259 & 267 = \texttt{BEO} \\
    12 & \texttt{OR}  & \texttt{T} & 261 & 268 = \texttt{ORT} \\
    13 & \texttt{TOB} & \texttt{E} & 266 & 269 = \texttt{TOBE} \\
    14 & \texttt{EO}  & \texttt{R} & 260 & 270 = \texttt{EOR} \\
    15 & \texttt{RN}  & \texttt{O} & 262 & 271 = \texttt{RNO} \\
    16 & \texttt{OT}  & (fin)      & 264 & --- \\
    \bottomrule
  \end{tabular}
\end{table}

\paragraph{El caso especial $KwKwK$.} Si la entrada contiene $K\,w\,K\,w\,K$, con $K$ un byte y
$w$ una cadena, el codificador puede emitir un código que acaba de crear y que el decodificador
aún no conoce. Con la entrada \texttt{aaa}, por ejemplo, el codificador emite 97
(\texttt{a}), agrega $257 = \texttt{aa}$ y luego emite 257. Al recibir 257, el decodificador
tiene $s = 257$: el código es igual al siguiente libre. La única frase consistente es la
anterior seguida de su propio primer byte, $\texttt{a} + \texttt{a} = \texttt{aa}$, que es lo
que hace la rama $k = s$ del \cref{alg:lzw-dec}; la prueba \texttt{handlesCodeNotYetInDictionary}
cubre este caso.

\paragraph{Análisis de complejidad.} Sea $D = 2^b$ el tamaño máximo del diccionario.
\begin{itemize}
  \item \emph{Codificación:} el diccionario se implementa como una tabla hash de
        direccionamiento abierto con sondeo lineal y hashing de Fibonacci (multiplicación por
        \texttt{0x9E3779B1}), indexada por la clave $(w, c)$. Su capacidad es $2D$, de modo que
        el factor de carga nunca supera $\alpha = 0.5$. Cada byte de entrada provoca una
        búsqueda; bajo la hipótesis de hashing uniforme, una búsqueda fallida requiere a lo
        sumo $1/(1-\alpha) = 2$ sondeos esperados~\cite{cormen2022}. Cada código emitido cuesta
        $O(1)$ adicional. Un reinicio limpia la tabla en $O(D)$, pero ocurre como mucho una vez
        cada $D - 257$ códigos, y cada código consume al menos un byte, de modo que su costo
        amortizado es $O(1)$ por byte. En total, $\Theta(n)$ esperado.
  \item \emph{Decodificación:} cada código se resuelve con accesos a arreglos
        (\texttt{prefix}, \texttt{suffix}, \texttt{firstByte}, \texttt{phraseLength}) y cada
        byte de salida se escribe exactamente una vez al recorrer la frase hacia atrás; el
        costo es $\Theta(n)$ en el peor caso, sin dependencia de una función hash.
  \item \emph{Espacio adicional:} $O(D)$, independiente de $n$. Para $b = 16$, el codificador
        usa dos arreglos de $2^{17}$ enteros y el decodificador cuatro arreglos de $2^{16}$
        entradas. La cabecera ocupa solo 7 bytes (magia, $b$ y longitud original).
\end{itemize}
\end{aporte}

% =====================================================================
\subsection{Comparación teórica}
\label{sec:comparacion-teorica}
\begin{aporte}{C}
La \cref{tab:comparacion-teorica} resume las diferencias estructurales entre ambos
algoritmos tal como se implementaron. La diferencia esencial no es asintótica ---ambos son
lineales--- sino de \emph{modelo}: Huffman es estático y de dos pasadas, y explota solo la
distribución de frecuencias de los bytes; LZW es adaptativo y de una pasada, y explota
frases repetidas, pero paga al menos 9 bits por código emitido.

\begin{table}[H]
  \centering
  \caption{Comparación teórica de las implementaciones.}
  \label{tab:comparacion-teorica}
  \small
  \renewcommand{\arraystretch}{1.15}
  \begin{tabularx}{\linewidth}{@{}>{\raggedright}p{3.4cm} >{\raggedright}X >{\raggedright\arraybackslash}X@{}}
    \toprule
    Aspecto & Huffman canónico & LZW (9--$b$ bits, \texttt{CLEAR}) \\
    \midrule
    Tiempo de codificación & $\Theta(n + \sigma\log\sigma) = \Theta(n)$ & $\Theta(n)$ esperado (tabla hash) \\
    Tiempo de decodificación & $\Theta(n\Lbar)$, con $\Lbar < H_0 + 1 \le 9$ & $\Theta(n)$ en el peor caso \\
    Espacio adicional & $O(\sigma)$ & $O(2^b)$, independiente de $n$ \\
    Modelo & Estático, dos pasadas; requiere transmitir el código & Adaptativo, una pasada; el diccionario no se transmite \\
    Redundancia explotada & Frecuencias de símbolos (orden cero) & Frases repetidas \\
    Garantía de tasa & $H_0 \le \Lbar < H_0 + 1$ & Asintótica para fuentes estacionarias~\cite{ziv1978}; sin cota para $n$ finito \\
    Cabecera & 262 bytes & 7 bytes \\
    \bottomrule
  \end{tabularx}
\end{table}
\end{aporte}

\subsection{Hipótesis}
\label{sec:hipotesis}
\begin{aporte}{C}
Del análisis anterior se derivan cuatro hipótesis que guían los experimentos de la
\cref{sec:metodologia}:
\begin{enumerate}[label=\textbf{H\arabic*.}, leftmargin=*]
  \item \label{hip:h1} Los codificadores y decodificadores de ambos algoritmos escalan
        linealmente con el tamaño de la entrada: la pendiente del ajuste log--log del tiempo
        frente a $n$ es cercana a 1.
  \item \label{hip:h2} Los bits por byte de Huffman se ubican en $[H_0, H_0 + 1)$ y, en fuentes
        i.i.d., quedan cerca de $H_0$.
  \item \label{hip:h3} LZW supera a Huffman en datos con frases repetidas (texto natural y
        sintético), pierde frente a él en datos i.i.d. y expande los datos aleatorios
        uniformes ($\bpb > 8$).
  \item \label{hip:h4} El tiempo de decodificación de Huffman crece con la longitud media del
        código (aproximadamente, con la entropía), mientras que el de LZW es prácticamente
        insensible a ella.
\end{enumerate}
\end{aporte}
